\documentclass[11pt,a4paper]{article}
\usepackage[margin=2.5cm]{geometry}
\usepackage{amsmath,amssymb,amsthm}
\usepackage[hidelinks]{hyperref}

\newtheorem{theorem}{Theorem}
\newtheorem{lemma}[theorem]{Lemma}
\newtheorem{corollary}[theorem]{Corollary}
\theoremstyle{remark}
\newtheorem{remark}[theorem]{Remark}

\newcommand{\R}{\mathbb{R}}

\title{Robustness of the weighted BPX factor for $M+\sqrt\beta K$\\[2pt]
\large Proof of Assumption~3.4 of the preconditioning paper (draft for checking)}
\author{Martin Stoll \and John W. Pearson}
\date{Working draft, November 2026}

\begin{document}
\maketitle

\noindent\emph{Status.} Steps marked \textbf{[proved]} are proved here
completely. Steps marked \textbf{[cited]} use classical multilevel results,
stated precisely as Lemmas~\ref{lem:jackson} and~\ref{lem:scs}; their exact
sources and constants still have to be checked. Notation: $\nu:=\sqrt\beta>0$
(the paper uses $\varepsilon$ for the solver accuracy).

\section{Setting}

Let $\Omega=(0,1)^d$, $h_\ell=2^{-\ell}$, and let $V_\ell\subset H^1_0(\Omega)$ be the $Q_1$
finite element space on the uniform grid of mesh size $h_\ell$, $\ell=1,\dots,L$. The
spaces are nested, $V_1\subset\dots\subset V_L=:V$. For $u\in V_\ell$ let $y\in\R^{n_\ell}$ be its
vector of nodal values. Define
\[
a(u,v)=(u,v)+\nu(\nabla u,\nabla v),\qquad t_\ell:=\nu h_\ell^{-2}=\nu4^\ell,\qquad \mu_\ell:=1+t_\ell .
\]
The matrix of $a$ in the nodal basis of $V_\ell$ is $A_\ell=M_\ell+\nu K_\ell$; on the finest
level $A_L=A_\beta$ of the paper.

\paragraph{BPX factor.} $P_\ell$ maps nodal values on level $\ell$ to nodal values on
level $L$ of the same function (the natural embedding $V_\ell\subset V_L$; $P_L=I$).
$\operatorname{diag}(A_\ell)=\delta_\ell I$ is constant on the uniform grid. With weights $\omega_\ell>0$,
\[
C=\sum_{\ell=1}^L\omega_\ell^2P_\ell P_\ell^T,\qquad \omega_\ell^2=\frac{\tilde\theta_\ell}{\delta_\ell}.
\]
The paper uses $\theta_\ell=\min(1,t_\ell)^s$ with $s\in[0,1]$. Here we use the
\emph{normalized} weights
\[
\tilde\theta_\ell=\frac{\theta_\ell}{\theta_L}=\Big(\frac{\min(1,t_\ell)}{\min(1,t_L)}\Big)^s .
\]
Since $\tilde\theta_\ell$ differs from $\theta_\ell$ by the factor $\theta_L^{-1}$, which does not depend on
$\ell$, $C$ only changes by a positive scalar. So $\kappa(CA_\beta)$ is the same for
both choices. The \emph{interval} $[c_0,c_1]$ below, however, is independent of
$\beta$ only for the normalized weights (relevant for the heat problem,
Corollary~\ref{cor:heat}).

\paragraph{Claim (Assumption 3.4).} For $s\in(0,1]$ there are $0<c_0\le c_1$,
depending only on $d$ and $s$, with $\operatorname{spec}(CA_\beta)\subset[c_0,c_1]$ for all $L\ge1$ and
$\nu>0$.

\section{Reduction to a decomposition problem}

\begin{lemma}[Lions; Xu--Zikatanov]\label{lem:lions}\textbf{[proved]}
Let $C=\sum_kR_kD_k^{-1}R_k^T$ with $D_k$ symmetric positive definite and
$\sum_k\operatorname{range}(R_k)=\R^n$. Then $C$ is SPD and, for every $x\in\R^n$,
\[
x^TC^{-1}x=\min\Big\{\sum_ky_k^TD_ky_k\ :\ \sum_kR_ky_k=x\Big\}.
\]
\end{lemma}
\begin{proof}
$C$ is SPD because $x^TCx=\sum_k|D_k^{-1/2}R_k^Tx|^2=0$ forces $R_k^Tx=0$ for all $k$,
hence $x\perp\sum_k\operatorname{range}R_k=\R^n$. For any decomposition $x=\sum_kR_ky_k$, by
Cauchy--Schwarz,
\[
x^TC^{-1}x=\sum_ky_k^TR_k^TC^{-1}x
\le\Big(\sum_ky_k^TD_ky_k\Big)^{1/2}\Big(\sum_kx^TC^{-1}R_kD_k^{-1}R_k^TC^{-1}x\Big)^{1/2},
\]
and the last sum equals $x^TC^{-1}CC^{-1}x=x^TC^{-1}x$. Hence $x^TC^{-1}x\le\sum_ky_k^TD_ky_k$.
Equality holds for $y_k=D_k^{-1}R_k^TC^{-1}x$, which is a decomposition of $x$.
\end{proof}

Apply this with $R_\ell=P_\ell$ and $D_\ell=\omega_\ell^{-2}I=(\delta_\ell/\tilde\theta_\ell)I$; $P_L=I$ makes
the ranges span. Let $x$ be the nodal vector of $u\in V$ and $y_\ell$ that of
$u_\ell\in V_\ell$; then $\sum_\ell P_\ell y_\ell=x$ means $\sum_\ell u_\ell=u$, and $x^TA_\beta x=a(u,u)$. So
\begin{equation}\label{eq:lions}
x^TC^{-1}x=\min_{u=\sum_\ell u_\ell,\ u_\ell\in V_\ell}\ \sum_\ell\frac{\delta_\ell}{\tilde\theta_\ell}|y_\ell|^2 .
\end{equation}
Since $\operatorname{spec}(CA_\beta)\subset[c_0,c_1]\iff c_1^{-1}A_\beta\preceq C^{-1}\preceq c_0^{-1}A_\beta$, the claim is
equivalent to
\begin{equation}\label{eq:twosided}
c_1^{-1}\,a(u,u)\ \le\ x^TC^{-1}x\ \le\ c_0^{-1}\,a(u,u)\qquad\text{for all }u\in V.
\end{equation}

\section{From nodal coefficients to $L^2$ norms}

\begin{lemma}\label{lem:nodal}\textbf{[proved]} For $v\in V_\ell$ with nodal vector $y$,
\[
\Big(\frac{h_\ell}3\Big)^d|y|^2<\|v\|^2<h_\ell^d|y|^2,\qquad
\delta_\ell=\Big(\frac{2h_\ell}3\Big)^d\big(1+3d\,t_\ell\big)\in\Big[\Big(\frac{2h_\ell}{3}\Big)^d\mu_\ell,\ 3d\Big(\frac{2h_\ell}{3}\Big)^d\mu_\ell\Big].
\]
\end{lemma}
\begin{proof}
$\|v\|^2=y^TM_\ell y$ with $M_\ell=M_1^{\otimes d}$, and the 1D mass matrix
$\frac{h}6\operatorname{tridiag}(1,4,1)$ has eigenvalues $\frac h6(4+2\cos\phi_j)\in(\frac h3,h)$. The
diagonal entries are $(M_\ell)_{ii}=(2h/3)^d$ and
$(K_\ell)_{ii}=d\cdot\frac2h(\frac{2h}3)^{d-1}$, so
$\delta_\ell=(\frac{2h}3)^d+\nu\frac{2d}h(\frac{2h}3)^{d-1}=(\frac{2h}3)^d(1+3d\nu h^{-2})$. Finally
$\mu\le1+3dt\le3d\mu$ for $t\ge0$ and $d\ge1$.
\end{proof}

Write $\lambda_\ell:=\mu_\ell/\tilde\theta_\ell$ and
\[
N(u):=\inf_{u=\sum_\ell u_\ell,\ u_\ell\in V_\ell}\ \sum_\ell\lambda_\ell\|u_\ell\|^2 .
\]
By Lemma~\ref{lem:nodal}, for each term
$(\frac23)^d\lambda_\ell\|u_\ell\|^2\le\frac{\delta_\ell}{\tilde\theta_\ell}|y_\ell|^2\le3d\,2^d\lambda_\ell\|u_\ell\|^2$
(the upper bound uses $h_\ell^d|y_\ell|^2<3^d\|u_\ell\|^2$, the lower bound $\|u_\ell\|^2<h_\ell^d|y_\ell|^2$).
Taking the infimum in~\eqref{eq:lions},
\begin{equation}\label{eq:NC}
\big(\tfrac23\big)^dN(u)\ \le\ x^TC^{-1}x\ \le\ 3d\,2^d\,N(u).
\end{equation}
It remains to show $K_1^{-1}a(u,u)\le N(u)\le K_0\,a(u,u)$ with $K_0,K_1$ independent of
$L$ and $\nu$. Then~\eqref{eq:twosided} holds with
\[
c_0=\frac1{3d\,2^dK_0},\qquad c_1=\big(\tfrac32\big)^dK_1,\qquad
\frac{c_1}{c_0}\le3d\cdot3^d\cdot K_0K_1 .
\]
The factor $3d\cdot3^d$ comes only from Lemma~\ref{lem:nodal}, i.e. from Jacobi
scaling and the nodal basis. It is the ``mass-matrix floor'' seen in the
numerics.

\section{Two regimes and the weights}

Since $t_\ell=\nu4^\ell$ increases by the factor 4 per level, there are two
regimes.
\begin{itemize}
\item \emph{Unresolved} ($t_L<1$): all levels are mass dominated. Then
  $\tilde\theta_\ell=(t_\ell/t_L)^s=4^{-s(L-\ell)}$ and $\lambda_\ell=(1+t_\ell)\,4^{s(L-\ell)}$.
\item \emph{Resolved} ($t_L\ge1$): let $\ell^*=\min\{\ell:t_\ell\ge1\}$. Then $\tilde\theta_\ell=\theta_\ell$, and
  $\lambda_\ell=1+t_\ell\in[t_\ell,2t_\ell]$ for $\ell\ge\ell^*$, while
  $\lambda_\ell=(1+t_\ell)t_\ell^{-s}\in[t_\ell^{-s},2t_\ell^{-s}]$ for $\ell<\ell^*$. Moreover $t_{\ell^*}<4$ if
  $\ell^*>1$.
\end{itemize}

\begin{lemma}\label{lem:lambda}\textbf{[proved]} For $s\in(0,1]$, in both regimes,
\[
\text{(a)}\quad t_\ell\le\lambda_\ell\ \text{ for all }\ell,\qquad
\text{(b)}\quad S:=\sum_{\ell=1}^L\lambda_\ell^{-1}\le\frac43+\frac1{1-4^{-s}} .
\]
For $s=0$, (a) still holds, but $S$ is of the order of the number of
mass-dominated levels: $S\approx L$ in the unresolved regime and $S\approx\ell^*\approx\log_4(1/\nu)$
in the resolved one.
\end{lemma}
\begin{proof}
(a) If $t_\ell\ge1$ then $\lambda_\ell=1+t_\ell>t_\ell$. If $t_\ell<1$ then $\lambda_\ell\ge1>t_\ell$, since
$\tilde\theta_\ell\le1$ in both regimes.
(b) Unresolved: $\lambda_\ell^{-1}\le4^{-s(L-\ell)}$, so $S\le\sum_{j\ge0}4^{-sj}=1/(1-4^{-s})$.
Resolved: for $\ell\ge\ell^*$, $\lambda_\ell^{-1}\le t_\ell^{-1}=t_{\ell^*}^{-1}4^{-(\ell-\ell^*)}\le4^{-(\ell-\ell^*)}$, which sums to
at most $4/3$. For $\ell<\ell^*$, $\lambda_\ell^{-1}\le t_\ell^s=t_{\ell^*-1}^s4^{-s(\ell^*-1-\ell)}<4^{-s(\ell^*-1-\ell)}$, which
sums to at most $1/(1-4^{-s})$. For $s=0$ the bounds $\lambda_\ell^{-1}\le1$ on the mass-dominated
levels are attained up to the factor $1+t_\ell<2$.
\end{proof}

\section{Two classical ingredients}

Let $Q_\ell:L^2(\Omega)\to V_\ell$ be the $L^2$-orthogonal projection, $Q_0:=0$.

\begin{lemma}[multilevel norm equivalence, Jackson direction]\label{lem:jackson}
\textbf{[cited]} There is $C_Q$, depending only on $d$, such that for all
$L$ and all $u\in V_L$
\[
\sum_{\ell=1}^Lh_\ell^{-2}\|(Q_\ell-Q_{\ell-1})u\|^2\le C_Q\,|u|_{H^1}^2 .
\]
\end{lemma}
\noindent Sources to check:
\begin{itemize}
\item Oswald, \emph{Multilevel Finite Element Approximation}, 1994;
\item Bornemann--Yserentant, Numer.\ Math.\ 64, 1993;
\item Dahmen--Kunoth, Numer.\ Math.\ 63, 1992;
\item Xu, SIAM Rev.\ 34, 1992.
\end{itemize}
Note that
the easy argument via $\|(I-Q_\ell)u\|\lesssim h_\ell|u|_1$ only gives $C_Q=O(L)$; the
$L$-independent bound is the deep part of BPX theory.

\begin{lemma}[strengthened Cauchy--Schwarz and inverse inequality]\label{lem:scs}
\textbf{[cited]} There are $C_S$ and $C_I$, depending only on $d$, such that
for $k\le\ell$, $u\in V_k$, $v\in V_\ell$,
\[
|(\nabla u,\nabla v)|\le C_S\,2^{-(\ell-k)/2}\,|u|_{H^1}\,h_\ell^{-1}\|v\|,\qquad
|v|_{H^1}^2\le C_Ih_\ell^{-2}\|v\|^2 .
\]
\end{lemma}
\noindent Sources to check: Yserentant (Numer.\ Math.\ 49, 1986), Xu (SIAM
Rev.\ 1992), Bramble--Pasciak--Xu (Math.\ Comp.\ 1990). For $Q_1$ on a uniform
grid the inverse inequality holds with $C_I=4d\cdot3^d$, by the eigenvalue bounds
$y^TK_\ell y\le4d\,h_\ell^{d-2}|y|^2$ and Lemma~\ref{lem:nodal}; an elementwise argument
gives a smaller constant. The strengthened Cauchy--Schwarz inequality uses
that $Q_1$ functions are harmonic on each coarse element (every $d$-linear
monomial is harmonic), integration by parts on coarse elements, and a
boundary-strip estimate for $v$.

\begin{lemma}[H$^1$ part of the upper bound]\label{lem:H1}\textbf{[proved from Lemma~\ref{lem:scs}]}
For every decomposition $u=\sum_\ell u_\ell$, $u_\ell\in V_\ell$,
\[
\nu|u|_{H^1}^2\le C_H\sum_\ell t_\ell\|u_\ell\|^2,\qquad C_H:=C_I+\frac{2C_SC_I^{1/2}}{\sqrt2-1}.
\]
\end{lemma}
\begin{proof}
Let $a_\ell:=t_\ell^{1/2}\|u_\ell\|=\nu^{1/2}h_\ell^{-1}\|u_\ell\|$. Expanding,
$\nu|u|_1^2=\nu\sum_\ell|u_\ell|_1^2+2\nu\sum_{k<\ell}(\nabla u_k,\nabla u_\ell)$. The diagonal is at most
$C_I\sum_\ell a_\ell^2$ by the inverse inequality. For $k<\ell$, Lemma~\ref{lem:scs} and the
inverse inequality for $u_k$ give
$\nu|(\nabla u_k,\nabla u_\ell)|\le C_SC_I^{1/2}2^{-(\ell-k)/2}a_ka_\ell$. With $q=2^{-1/2}$,
$\sum_{k<\ell}q^{\ell-k}a_ka_\ell\le\big(\sum_{j\ge1}q^j\big)\sum_\ell a_\ell^2=\frac1{\sqrt2-1}\sum_\ell a_\ell^2$ by the
Schur test ($2a_ka_\ell\le a_k^2+a_\ell^2$).
\end{proof}

\section{The two bounds}

\begin{lemma}[upper bound for $a$]\label{lem:upper}\textbf{[proved, given Lemma~\ref{lem:scs}]}
For $s\in(0,1]$ and every decomposition $u=\sum_\ell u_\ell$,
\[
a(u,u)\le K_1\sum_\ell\lambda_\ell\|u_\ell\|^2,\qquad K_1:=\frac43+\frac1{1-4^{-s}}+C_H .
\]
Consequently $a(u,u)\le K_1N(u)$.
\end{lemma}
\begin{proof}
Mass part: by Cauchy--Schwarz,
$\|\sum_\ell u_\ell\|\le\sum_\ell\lambda_\ell^{-1/2}\cdot\lambda_\ell^{1/2}\|u_\ell\|\le S^{1/2}(\sum_\ell\lambda_\ell\|u_\ell\|^2)^{1/2}$, and $S$ is
bounded by Lemma~\ref{lem:lambda}(b). Stiffness part: Lemma~\ref{lem:H1} and
$t_\ell\le\lambda_\ell$ (Lemma~\ref{lem:lambda}(a)).
\end{proof}

\begin{lemma}[stable decomposition]\label{lem:lower}\textbf{[proved, given Lemma~\ref{lem:jackson}]}
For every $u\in V$, $N(u)\le K_0\,a(u,u)$ with $K_0:=\max(5,2C_Q)$, for every
$s\in[0,1]$.
\end{lemma}
\begin{proof}
\emph{Unresolved regime.} Take $u_L=u$ and $u_\ell=0$ otherwise. Then $\lambda_L=1+t_L<2$
and $N(u)\le2\|u\|^2\le2a(u,u)$.

\emph{Resolved regime.} Take $u_\ell=0$ for $\ell<\ell^*$, $u_{\ell^*}=Q_{\ell^*}u$, and
$u_\ell=(Q_\ell-Q_{\ell-1})u$ for $\ell>\ell^*$. These sum to $u$, since $Q_L u=u$. For $\ell>\ell^*$,
$\lambda_\ell\le2t_\ell$. For $\ell^*$:
\begin{itemize}
\item if $\ell^*>1$: $\lambda_{\ell^*}=1+t_{\ell^*}<5$ and $\|Q_{\ell^*}u\|\le\|u\|$, so this term is at most
  $5\|u\|^2$;
\item if $\ell^*=1$: $\lambda_1\|Q_1u\|^2\le\|u\|^2+t_1\|(Q_1-Q_0)u\|^2$.
\end{itemize}
In both cases
\[
N(u)\le5\|u\|^2+2\nu\sum_{\ell\ge1}h_\ell^{-2}\|(Q_\ell-Q_{\ell-1})u\|^2\le5\|u\|^2+2C_Q\,\nu|u|_1^2\le K_0\,a(u,u)
\]
by Lemma~\ref{lem:jackson}.
\end{proof}

\section{Result}

\begin{theorem}\label{thm:robust}
Let $s\in(0,1]$ and use the normalized weights $\tilde\theta_\ell$. For all $L\ge1$ and
$\beta>0$,
\[
\operatorname{spec}(CA_\beta)\subset[c_0,c_1],\qquad c_0=\frac1{3d\,2^d\max(5,2C_Q)},\qquad
c_1=\big(\tfrac32\big)^d\Big(\frac43+\frac1{1-4^{-s}}+C_H\Big).
\]
In particular $\kappa(CA_\beta)\le3d\cdot3^d\max(5,2C_Q)\big(\frac43+\frac1{1-4^{-s}}+C_H\big)$,
independently of $h$ and $\beta$. The same condition number bound holds
for the unnormalized weights $\theta_\ell$ of the paper.
\end{theorem}
\begin{proof}
Combine~\eqref{eq:NC} with Lemmas~\ref{lem:upper} and~\ref{lem:lower}. The
last sentence follows from the scale invariance of $\kappa$.
\end{proof}

\begin{remark}[$s=0$, plain BPX]
Lemma~\ref{lem:lower} still holds, but in Lemma~\ref{lem:upper} the mass
part contributes $S$, which is of the order of the number of mass-dominated
levels. The resulting bound grows like $\log(1/\beta)$ in the resolved regime
and like $L$ in the unresolved regime. This matches the observed
degradation of plain BPX. It is also consistent with the $L$-dependence in
Yang--Yu--Zhang (arXiv:2610.06987), whose weights do not damp mass-dominated
levels (to be checked against their proof).
\end{remark}

\begin{remark}[sharpness]
The constants are far from sharp: for $d=3$ the factor $3d\cdot3^d=243$ alone
exceeds the observed $\kappa(CA_\beta)\le20.5$. The proof explains the structure
(geometric damping makes the mass part summable; the floor comes from
Jacobi scaling), not the size of the constants.
\end{remark}

\begin{corollary}[time-periodic heat problem]\label{cor:heat}
For frequency $k$ the paper's operator is $A_k=r_kM+\nu K=r_k(M+\nu_kK)$ with
$\nu_k=\nu/r_k$. Using in $R_k$ the diagonal of $A_k$ and the normalized weights
for $\nu_k$, we get $C_kA_k=C(\nu_k)A_{\nu_k}$ with $C(\nu_k)$ the factor of
Theorem~\ref{thm:robust} for parameter $\nu_k$. Hence
$\operatorname{spec}(C_kA_k)\subset[c_0,c_1]$ with the \emph{same} $c_0,c_1$ for all $k$. This is the
hypothesis of the paper's heat theorem.
\end{corollary}
\begin{proof}
$\operatorname{diag}(A_{k,\ell})=r_k\operatorname{diag}(M_\ell+\nu_kK_\ell)$, so $\omega_{\ell,k}^2=r_k^{-1}\omega_\ell(\nu_k)^2$ and
$C_k=r_k^{-1}C(\nu_k)$. Then $C_kA_k=C(\nu_k)(M+\nu_kK)$, and Theorem~\ref{thm:robust}
applies with $\nu_k$; its interval does not depend on $\nu$.
\end{proof}

\end{document}
